% 完备空间（综述）
% license CCBYSA3
% type Wiki

本文根据 CC-BY-SA 协议转载翻译自维基百科\href{https://en.wikipedia.org/wiki/Complete_metric_space}{相关文章}

在数学分析中，如果一个度量空间$M$中的每个柯西序列都收敛于$M$中的某个点，那么这个度量空间$M$就称为完备的，也称为柯西空间。

直观地说，一个空间是完备的，意味着它内部或边界上不存在“缺失的点”。例如，有理数集并不是完备的，因为$\sqrt{2}$这样的数“缺失”在有理数集中，尽管我们可以构造一个收敛于$\sqrt{2}$的有理数柯西序列（见下文更多示例）。总是可以通过“填补所有这些空缺”来得到某个空间的完备化，具体方法将在下文解释。
\subsection{定义}
\textbf{柯西序列}

一个度量空间$(X, d)$中的序列$x_{1}, x_{2}, x_{3}, \ldots$称为柯西序列，如果对于任意一个正实数$r > 0$，存在一个正整数$N$，使得对于所有大于$N$的正整数 $m, n$，都有：
$$
d(x_{m}, x_{n}) < r.~
$$
\textbf{完备空间}

一个度量空间$(X, d)$称为完备的，当且仅当满足以下任一等价条件：
\begin{enumerate}
\item 柯西序列收敛性，$X$ 中的每一个柯西序列都在 $X$ 中收敛（即收敛到 $X$ 中的某个点）。
\item 嵌套闭集的交集不为空，$X$ 中任意一个非空闭集序列$\{F_n\}$，如果满足以下条件：$F_{n+1} \subseteq F_n$（即序列是逐步缩小的嵌套闭集序列）；集合的直径 $\operatorname{diam}(F_n)$ 趋于 0：$\operatorname{diam}(F_n) \to 0$，那么这些集合的交集$\bigcap_{n=1}^{\infty} F_n$恰好包含一个唯一的点 $x \in X$。
\end{enumerate}
\subsection{示例}
有理数集 $\mathbb{Q}$ 搭配标准度量（即差的绝对值）并不是完备的。例如，考虑如下定义的序列：
$$
x_1 = 1, \quad x_{n+1} = \frac{x_n}{2} + \frac{1}{x_n}.~
$$
这是一个由有理数组成的柯西序列，但它并不收敛于任何一个有理数。如果这个序列确实有极限 $x$，那么根据递推公式有：$x = \frac{x}{2} + \frac{1}{x}$，从而得到$
x^2 = 2$。然而，没有任何一个有理数满足这个条件。但如果把该序列放在实数集里来看，它确实收敛到无理数 $\sqrt{2}$。

开区间 $(0,1)$ 也不是完备的（度量依然是绝对值差）。考虑序列：$x_n = \frac{1}{n}$。这是一个柯西序列，但在 $(0,1)$ 中没有极限。然而，闭区间 $[0,1]$ 是完备的；同样的序列在这个空间中有极限，其极限点就是 $0$。

实数集 $\mathbb{R}$ 和复数集 $\mathbb{C}$（度量由差的绝对值给出）都是完备的。同样，带有通常欧几里得距离度量的欧几里得空间 $\mathbb{R}^n$ 也是完备空间。相比之下，无限维的赋范向量空间可能是完备的，也可能不是。那些完备的空间称为巴拿赫空间。例如，闭有界区间$[a, b]$上的实值连续函数空间 $C[a,b]$，在上确界范数下是一个巴拿赫空间，因此也是一个完备的度量空间。然而，对于开区间 $(a, b)$ 上的连续函数空间 $C(a,b)$，上确界范数并不能定义一个范数，因为该空间中可能存在无界函数。相反，如果赋予其紧致收敛拓扑，则 $C(a,b)$ 可以被赋予弗雷歇空间的结构，即一个局部凸的拓扑向量空间，其拓扑可以由一个平移不变的完备度量诱导。


对于任意素数 $p$，$p$ 进数空间 $\mathbf{Q}_p$ 是完备的。该空间通过 $p$ 进度量将有理数集 $\mathbf{Q}$ 完备化，正如实数集 $\mathbf{R}$ 通过通常度量将 $\mathbf{Q}$ 完备化一样。

如果 $S$ 是一个任意集合，那么 $S^{\mathbf{N}}$（即 $S$ 中所有序列的集合）在如下定义的距离下也是一个完备度量空间：对任意两个序列 $\left(x_n\right)$ 和 $\left(y_n\right)$，它们的距离定义为$\frac{1}{N}$，其中 $N$ 是第一个满足 $x_N \neq y_N$ 的索引；如果没有这样的索引（即两个序列完全相同），则距离为 $0$。该空间与可数个离散空间$S$的乘积空间同胚。

完备的黎曼流形称为测地流形。这种完备性可以由Hopf–Rinow 定理推出。
\subsection{一些定理}
紧致空间与完备性,每个紧致的度量空间都是完备的，但完备空间不一定是紧致的。事实上，一个度量空间当且仅当它是完备的且全有界的时才是紧致的。这推广了Heine–Borel 定理，该定理指出 $\mathbb{R}^n$ 中任意闭且有界的子空间$S$是紧致的，因此也是完备的。\(^\text{[1]}\)

闭集与完备性,若$(X,d)$是一个完备的度量空间，且$A \subseteq X$是闭集，则$A$也是完备的。反过来，如果$(X,d)$是一个度量空间，且$A \subseteq X$是一个完备的子空间，则$A$必然是闭集。

有界函数空间的完备性,如果$X$是一个集合，且$M$是一个完备的度量空间，则所有从 $X$到$M$的有界函数所组成的集合$B(X,M)$是一个完备的度量空间。这里，$B(X,M)$ 中的距离通过 $M$ 中的度量定义，并采用上确界范数：
  $$
  d(f,g) \equiv \sup \{ d[f(x),g(x)] : x \in X \}.~
  $$
连续有界函数空间的完备性,如果$X$是一个拓扑空间，且$M$是一个完备的度量空间，那么所有从$X$到$M$的连续有界函数构成的空间$C_b(X,M)$是$B(X,M)$的闭子空间，因此它也是完备的。

Baire 类定理,Baire 类定理指出，每个完备的度量空间都是一个Baire 空间。也就是说，这个空间中任意可数多个稠密度为零的子集的并集都没有内部点。

Banach 不动点定理,Banach 不动点定理指出，在一个完备的度量空间中，任何一个压缩映射都存在一个不动点。该定理常用于证明完备度量空间（如巴拿赫空间）上的反函数定理。

定理[2]（C. Ursescu）——设$X$是一个完备的度量空间，且$S_1, S_2, S_3, \ldots$是一列 $X$ 的子集。
\begin{itemize}
\item 当每个 $S_i$ 都是闭集时：$\operatorname{cl} \left(\bigcup_{i\in\mathbb{N}} \operatorname{int} S_i\right) = \operatorname{cl}\operatorname{int} \left(\bigcup_{i\in\mathbb{N}} S_i\right)$其中 $\operatorname{cl}$ 表示闭包，$\operatorname{int}$ 表示内部。
\item 当每个$S_i$都是开集时：$\operatorname{int} \left(\bigcap_{i\in\mathbb{N}} \operatorname{cl} S_i\right)=\operatorname{int} \operatorname{cl} \left(\bigcap_{i\in\mathbb{N}} S_i\right)$
\end{itemize}
\subsection{完备化}
对于任意一个度量空间 $M$，都可以构造一个完备的度量空间$M'$（也记作$\overline{M}$），其中 $M$ 是 $M'$ 的一个稠密子空间。该完备空间具有如下泛性质：如果 $N$ 是任意一个完备的度量空间，且 $f$ 是从 $M$ 到 $N$ 的任意一致连续函数，那么存在唯一的一个从 $M'$ 到 $N$ 的一致连续函数 $f'$，它是 $f$ 的延拓。满足这一性质的 $M'$（在所有同构地包含 $M$ 的完备度量空间中）是唯一确定的，这个空间称为$M$ 的完备化。

$M$ 的完备化可以通过柯西序列的等价类来构造。设 $M$ 中的两列柯西序列分别为$x_{\bullet} = (x_n) \quad \text{和} \quad y_{\bullet} = (y_n)$,定义它们之间的“距离”为：
$$
d(x_{\bullet}, y_{\bullet}) = \lim_{n} d(x_n, y_n)~
$$
这个极限一定存在，因为实数集是完备的。不过，此时定义的只是一个伪度量，还不是严格意义上的度量，因为可能存在不同的柯西序列之间的距离为 0 的情况。但“距离为 0”这个关系在所有柯西序列上定义了一个等价关系。将这些柯西序列按该关系分类得到的等价类集合，就构成了一个真正的度量空间，这就是$M$的完备化。原始空间$M$可以自然嵌入到$M'$中：将$M$中的一个元素$x$对应于那个常值序列（即所有项都是$x$的序列）的等价类。这样就得到一个到稠密子空间的同构等距映射，满足完备化的定义要求。需要注意的是，这种构造明确利用了实数集的完备性，因此在对有理数进行完备化时，需要稍作不同的处理。

康托尔对实数的构造方法与上述完备化过程类似：实数集可以看作是有理数集在普通绝对值度量下的完备化。不过，这里有一个微妙的问题：在构造实数时，逻辑上不允许直接依赖实数本身的完备性。尽管如此，依然可以按相同的方法定义有理数柯西序列的等价类，并证明这些等价类构成的集合是一个域，并且该域包含有理数作为其子域。这个域是完备的，且自然带有全序结构，并且是唯一的（在同构意义下）完全有序完备域。这个域就是实数域（详细内容可参见“实数的构造”）。

从直观上理解，每一个“应该”收敛到某个实数的有理数柯西序列所形成的等价类，就对应于那个实数。以十进制展开的截断序列，只是该等价类中的一个典型柯西序列。

对于任意素数 $p$，$p$进数是通过在不同的度量下对有理数集进行完备化得到的。

如果将上述完备化过程应用于一个赋范向量空间，就得到一个包含原空间作为稠密子空间的巴拿赫空间。同样，如果将这一过程应用于一个内积空间，则结果是一个包含原空间作为稠密子空间的希尔伯特空间。
\subsection{拓扑完备空间}
完备性是度量的性质，而不是拓扑的性质。这意味着，一个完备的度量空间可以与一个非完备的空间同胚。例如，实数集$\mathbb{R}$是完备的，但它与开区间$(0,1)$同胚，而后者并不是完备的。

在拓扑学中，人们研究完全可度量化空间，即存在至少一个完备度量能够诱导其给定拓扑的空间。完全可度量化空间可以刻画为：某个完备度量空间中可数多个开子集的交集。由于 Baire 类定理的结论是纯拓扑性的，因此它同样适用于这些空间。

完全可度量化空间常被称为拓扑完备空间。然而，“拓扑完备”这一称谓并非完全精确，因为度量并不是讨论完备性的拓扑空间中最一般的结构（参见“替代与推广”部分）。事实上，有些作者将“拓扑完备”用于更广泛的一类拓扑空间，即完全可一致化空间。\(^\text{[2]}\)

与某个可分且完备的度量空间同胚的拓扑空间称为波兰空间。
\subsection{替代与推广}
由于柯西序列也可以在一般的拓扑群中定义，因此除了依赖度量结构来定义完备性和构造空间的完备化之外，还可以利用群结构来进行推广。这种方法常见于拓扑向量空间的研究中，但它只需要存在一个连续的“减法”运算即可。在这种情境下，比较两点$x$和$y$的“距离”不再通过度量$d$和实数$\varepsilon$的不等式$d(x, y)<\varepsilon$来衡量，而是通过减法比较它们的差是否属于某个$0$的开邻域 $N$：$x - y \in N$.

这些定义的一个常见推广出现在一致空间的框架下，在一致空间中，“邻域”是指所有距离不超过某个“范围”的点对的集合。

在定义完备性时，还可以将柯西序列替换为柯西网或柯西滤子。如果空间$X$中的每个柯西网（或等价地，每个柯西滤子）都有极限，那么称$X$是完备的。此外，可以像度量空间的完备化那样，为任意一致空间构造它的完备化。在更一般的情形下，柯西空间也可以使用柯西网的概念，它们同样有类似于一致空间的完备性和完备化的概念。
\subsection{参见}
\begin{itemize}
\item 柯西空间 —— 广义拓扑与分析中的概念
\item 完备化（代数） —— 代数中关于理想幂的完备化
\item 完备一致空间 —— 具备一致性结构的拓扑空间
\item 完备拓扑向量空间—— 泛函分析中的一种结构
\item 埃克兰变分原理
\item Knaster–Tarski 定理 —— 序理论与格理论中的一个定理
\end{itemize}
\subsection{注释}
\begin{enumerate}
\item Sutherland, Wilson A. (1975). Introduction to Metric and Topological Spaces. Clarendon Press. ISBN 978-0-19-853161-6.
\item Zalinescu, C. (2002). Convex analysis in general vector spaces. River Edge, N.J. London: World Scientific. p. 33. ISBN 981-238-067-1. OCLC 285163112.
\item Kelley, Problem 6.L, p. 208
\end{enumerate}
\subsection{参考文献}
\begin{itemize}
\item Kelley, John L. (1975). General Topology. Springer. ISBN 0-387-90125-6.
\item Kreyszig, Erwin. Introductory functional analysis with applications. Wiley, New York, 1978. ISBN 0-471-03729-X.
\item Lang, Serge. Real and Functional Analysis. ISBN 0-387-94001-4.
\item Meise, Reinhold; Vogt, Dietmar (1997). Introduction to functional analysis. Ramanujan, M.S. (译). Oxford: Clarendon Press; New York: Oxford University Press. ISBN 0-19-851485-9.
\end{itemize}
